\documentclass{report}
\usepackage{amsmath,amssymb}
\setlength{\parindent}{0mm}
\setlength{\parskip}{1em}
\begin{document}
\begin{center}
\rule{6in}{1pt} \
{\large Alexander, J.M. Howse \\
{\bf Nonlinearly Preconditioned Optimization on Grassmann Manifolds for Tucker Tensor Approximations}}

University of Waterloo \\ Department of Applied Mathematics \\ 200 University Ave W \\ Waterloo \\ ON N2L 3G1 \\ Canada
\\
{\tt ajmhowse@gmail.com}\\
Hans De Sterck\end{center}

Two accelerated optimization algorithms are presented for computing
approximate Tucker tensor decompositions (ATDs). The first is a
nonlinearly preconditioned conjugate gradient (NPCG) algorithm, wherein a
nonlinear preconditioner generates a direction replacing the gradient in
the nonlinear conjugate gradient iteration. The second is a nonlinear
GMRES (N-GMRES) algorithm, in which a linear combination of iterates
generated by a nonlinear preconditioner is minimized to produce an
improved search direction. The Euclidean versions of these methods are
extended to the manifold setting, where optimization on Grassmann
manifolds is used to handle orthonormality constraints and to allow
isolated minimizers. The higher order orthogonal iteration (HOOI), the
workhorse algorithm for computing ATDs, is used as the nonlinear
preconditioner in NPCG and N-GMRES. Four options are provided for the
update parameter $\beta$ in NPCG. Two strategies for approximating the
Hessian operator applied to a vector are provided for N-GMRES. NPCG and
N-GMRES are compared to HOOI, NCG, limited memory BFGS, and a manifold
trust region algorithm using synthetic data and real life tensor data
arising from handwritten digit recognition. Numerical results show that
all four NPCG variants and N-GMRES using a difference of gradients
Hessian approximation accelerate HOOI significantly for large tensors,
noisy data, and when high accuracy results are required. For these
problems, the proposed methods converge faster and more robustly than
HOOI and the state-of-the-art methods considered.


\end{document}
